\section{La fuente s\'ismica: el martillo, sus puntas y el martinete de leva}
\label{sec:fuente-sismica}

La fuente es la mitad del experimento y es, con diferencia, la parte menos
controlada del sistema. Esta secci\'on re\'une lo que se sabe de ella, incluyendo
un dato del manual del instrumento que estuvo seis meses sin explotarse y que
resulta directamente accionable.

\subsection{El instrumento: PCB Piezotronics 086D20}

\pendiente{\textbf{Correcci\'on de terminolog\'ia.} Durante meses los documentos
del proyecto describieron el 086D20 como un ``aceler\'ometro''. \textbf{No lo
es}: seg\'un su propio manual es un \emph{Short-sledge Impulse Hammer with force
sensor \& tips}, es decir un martillo con \textbf{sensor de fuerza} de cuarzo
ICP\textsuperscript{\textregistered} montado en la cara de impacto. La bit\'acora
del 5 de mayo lo registr\'o correctamente como ``sensor de fuerza
piezoel\'ectrico''; el error se introdujo despu\'es, al redactar. Corregido el 12
de agosto en los cinco lugares donde aparec\'ia.

La distinci\'on importa: el canal de martillo mide \textbf{fuerza aplicada al
suelo en newtons}, no aceleraci\'on de la maza. Eso hace que el canal sea
directamente interpretable como la funci\'on de fuente del experimento ---y, en
principio, utilizable para deconvoluci\'on--- cosa que un aceler\'ometro sobre la
maza no permitir\'ia sin un modelo de la masa efectiva.}

\begin{center}\footnotesize
\begin{tabularx}{\textwidth}{@{}l X@{}}
\toprule
Par\'ametro & Valor \\
\midrule
Tipo & \emph{Short-sledge impulse hammer} con sensor de fuerza integrado \\
Rango de fuerza & 0 a \num{5000}~lbf \\
Sensibilidad & \SI{1}{\milli\volt}/lbf $=$ \SI{0.23}{\milli\volt\per\newton} \\
Masa de la cabeza & \num{2.4}~lb ($\approx\SI{1.09}{\kilo\gram}$) \\
Tipo de sensor & Cuarzo ICP\textsuperscript{\textregistered}, amplificador de ganancia unitaria integrado \\
Alimentaci\'on & Dos hilos, corriente constante; salida acoplada en alterna \\
Construcci\'on & Unidad integral soldada por l\'aser; \textbf{no se desarma fuera de f\'abrica} \\
Cable & Cinta reparable, dise\~nada a prop\'osito como el eslab\'on d\'ebil ante un golpe fallido \\
\bottomrule
\end{tabularx}
\end{center}

\subsection{Las puntas: el control de ancho de banda de la fuente}

\dato{\textbf{El hallazgo}. El manual incluye las curvas de respuesta en
frecuencia de la familia 086D20 para cada punta de impacto, todas a
\num{1000}~lbf nominales. La punta \emph{no} es un accesorio de
protecci\'on: \textbf{es el \'unico control de ancho de banda de la fuente que
tiene el operador}. Valores le\'idos de las curvas (ca\'ida de
\SI{-20}{\decibel}):

\begin{center}
\begin{tabular}{@{}ll@{}}
\toprule
Punta & \SI{-20}{\decibel} aproximadamente en \\
\midrule
Super blanda & \SI{600}{\hertz} \\
Blanda & \SI{700}{\hertz} \\
Media & \SI{900}{\hertz} \\
Media/dura & \SI{1300}{\hertz} \\
\bottomrule
\end{tabular}
\end{center}

Cuanto m\'as blanda la punta, m\'as largo el pulso de fuerza, m\'as estrecho su
espectro y \textbf{m\'as concentrada la energ\'ia en baja frecuencia}. El manual
lo dice sin rodeos: ``las puntas de distinta rigidez permiten variar el ancho del
pulso y el contenido en frecuencia de la fuerza''.}

\seleccionado{\textbf{Consecuencia directa para este proyecto, y es fuerte.}
Todo el trabajo de electr\'onica persigue el extremo inferior de la banda porque
ah\'i est\'a la informaci\'on de profundidad (\S\ref{sec:atlas} y el argumento
$\lambda$--$f$). Pero la fuente tiene su propio control de banda, gratuito y del
lado correcto: \textbf{usar la punta m\'as blanda posible concentra la energ\'ia
del golpe exactamente en la banda que el sistema quiere medir}. Es la mejora de
relaci\'on se\~nal--ruido de menor costo disponible ---no cuesta un componente ni
una l\'inea de c\'odigo, s\'olo enroscar otra punta.

El manual agrega un segundo recurso en la misma direcci\'on: la \emph{masa de
extensi\'on} que se provee con la mayor\'ia de los martillos ``permite un ajuste
adicional concentrando m\'as energ\'ia en frecuencias bajas''. Ambos recursos
apuntan al mismo lado y ninguno de los dos se explot\'o de forma sistem\'atica en
las campa\~nas realizadas.}

\pendiente{\textbf{Esto explica retroactivamente una decisi\'on de interfaz.} La
aplicaci\'on \texttt{geophone\_scope\_simple.m} tiene un selector de punta de
martillo con las opciones \emph{gris}, \emph{roja}, \emph{marr\'on}, \emph{negra}
y \emph{mezclado}, y ese dato se guarda con cada ensayo. En su momento parec\'ia
mera prolijidad de registro. A la luz de las curvas del manual es otra cosa:
\textbf{es el registro de la banda de la fuente en cada captura}. Queda
pendiente mapear cada color a su dureza y contrastar el contenido espectral
medido de las capturas con las curvas del fabricante ---un experimento que se
puede hacer \'integramente con datos ya adquiridos.}

\subsection{El martinete de leva: hacia una fuente repetible}

El problema del martillo manual no es su banda sino su \textbf{repetibilidad}:
la energ\'ia, el \'angulo y el punto exacto del golpe cambian entre repeticiones,
y el m\'etodo de apertura sint\'etica adoptado (mover un ge\'ofono y apilar golpes
sucesivos) \emph{depende} de que la fuente sea repetible ---es su hip\'otesis
m\'as fuerte y la menos verificada.

De ah\'i el dise\~no de un \textbf{martinete de leva}: un mecanismo que eleva una
maza siguiendo un perfil de leva y la libera siempre desde la misma altura, con
la misma energ\'ia. Se model\'o en Simscape Multibody a partir de un CAD de
Fusion~360.

\begin{center}\footnotesize
\begin{tabular}{@{}lr@{}}
\toprule
Cota (del CAD, \texttt{lista\_cotas\_herrero.csv}) & Valor \\
\midrule
Pivote a rodillo & \SI{215}{\milli\meter} \\
\'Angulo de palanca & \SI{127.5}{\degree} \\
Palanca: longitud $\times$ ancho $\times$ espesor & $195\times24\times\SI{24}{\milli\meter}$ \\
Eje del rodillo: di\'ametro / longitud & \SI{16}{} / \SI{72}{\milli\meter} \\
Leva: radio m\'inimo / m\'aximo & \SI{5.329}{\centi\meter} / \SI{17.760}{\centi\meter} \\
Separaci\'on entre planos & \SI{60}{\milli\meter} \\
Holgura axial leva--palanca & \SI{39}{\milli\meter} \\
\bottomrule
\end{tabular}
\end{center}

\begin{figure}[ht]
\centering
\includegraphics[width=\textwidth]{mosaico_martillo_leva.png}
\caption{Modelado del martinete de leva: perfil de leva a \SI{32}{\degree},
plano de taller en A3, y los resultados de simulaci\'on multibody ---caso
nominal, subida, multivuelta contra el piso y la transici\'on de toma C2.}
\label{fig:mosaico-martillo}
\end{figure}

La cadena de herramientas qued\'o completa y trazable:

\begin{center}\footnotesize
\begin{tabularx}{\textwidth}{@{}l X@{}}
\toprule
\texttt{extraer\_geometria\_cad.m} & Mide el mecanismo directamente del CAD de Fusion 360 y lo cachea en \texttt{geometria\_cad.mat} \\
\texttt{params\_martillo.m} & Separa lo medido de lo decidido: material (PLA, PETG, ABS, acero), rpm de la leva, rigideces de contacto y el factor de escala \\
\texttt{construir\_martillo\_leva.m} & Arma el modelo Simscape Multibody \\
\texttt{contacto\_leva\_seguidor.m} & Modelo de contacto entre leva y seguidor \\
\texttt{simular\_martillo.m} & Corre la simulaci\'on y reporta resultados \\
\bottomrule
\end{tabularx}
\end{center}

\dato{\textbf{El factor de escala es una decisi\'on de dise\~no expl\'icita.} El
CAD est\'a dibujado para \textbf{imprimir en 3D}: las formas son las definitivas
pero el tama\~no es de maqueta. El par\'ametro \texttt{escala}${}=10$ recupera el
dise\~no de campo, con $C=(250,400)$~mm, que es exactamente diez veces el
$(25,40)$~mm del CAD. Es decir: se valida la cinem\'atica en una maqueta
imprimible y barata, y el mismo modelo escala al mecanismo real sin rehacer nada.
Los par\'ametros por defecto son \SI{40}{rpm} y PLA
($\rho=\SI{1240}{\kilo\gram\per\meter\cubed}$).}

\seleccionado{Existe un \textbf{plano de taller en A3} listo para entregar a un
herrero (\texttt{plano\_herrero\_martinete\_A3.pdf/.svg/.png}), con la lista de
cotas en CSV y el perfil de leva exportado en DXF a escala 1:1
(\texttt{perfil\_leva\_1a1\_mm.dxf}) y en STL para impresi\'on
(\texttt{leva\_balistica\_12lb.stl}). El trabajo lleg\'o hasta el punto en que
\emph{s\'olo falta fabricarlo}.}

\pendiente{El martinete no se construy\'o ni se ensay\'o. Mientras tanto, la
hip\'otesis de repetibilidad de la fuente ---sobre la que descansa toda la
apertura sint\'etica--- segu\'ia sin verificarse. La secci\'on siguiente la
verifica con datos que ya exist\'ian.}

\subsection{Repetibilidad medida de la fuente (an\'alisis del 12/08/2026)}
\label{sec:repetibilidad-fuente}

Este an\'alisis es nuevo: se hizo al escribir este historial, sobre los
\textbf{646 golpes} del conjunto procesado \texttt{Canchita}, repartidos en
\textbf{22 posiciones}. Para cada golpe se midi\'o, en la ventana del impacto del
canal de martillo, la amplitud pico, el ancho del pulso a media altura y el
centroide espectral. Los datos estaban desde julio; el c\'alculo nunca se hab\'ia
hecho.

\begin{figure}[ht]
\centering
\includegraphics[width=\textwidth]{repetibilidad_fuente.png}
\caption{Repetibilidad del golpe medida sobre el canal de martillo. Puntos
claros: cada golpe individual. L\'ineas: mediana por posici\'on. La banda gris
del panel superior marca el hueco vac\'io que separa las dos configuraciones.}
\label{fig:repetibilidad-fuente}
\end{figure}

\dato{\textbf{Hallazgo 1: el conjunto procesado mezcla dos configuraciones de
adquisici\'on.} La distribuci\'on de amplitudes es netamente \textbf{bimodal},
con un \textbf{hueco absolutamente vac\'io entre \SI{3.08}{} y
\SI{3.22}{\volt}} ---ni un solo golpe de los 646 cae ah\'i. Los
\textbf{191 golpes tomados a \SI{2929}{\hertz} caen todos} del lado alto. La
nomenclatura de las carpetas confirma que son sesiones distintas
(\texttt{muestra20} frente a \texttt{muestras10},
\texttt{muestra18metrosCapaz}, \texttt{muestra16metroslegit}: los sufijos
``legit'' y ``Capaz'' delatan decisiones tomadas sobre la marcha en el campo).
\textbf{Promediar las dos poblaciones juntas fabrica una dispersi\'on que no
existe}: hacerlo daba un CV de amplitud del \SI{20}{\percent}, que es un
artefacto de la mezcla y no una propiedad del martillo.}

\dato{\textbf{Hallazgo 2: en la configuraci\'on A el canal de martillo est\'a
saturado.} Separadas las poblaciones, la configuraci\'on A da un CV de amplitud
de \textbf{\SI{0.4}{\percent}} mediano. Eso es f\'isicamente imposible para un
martillo balanceado a mano. Desglosado por sesi\'on, el cuadro es inequ\'ivoco:

\begin{center}
\begin{tabular}{@{}lrrr@{}}
\toprule
Sesi\'on & $n$ & Pico mediano [\si{\volt}] & CV [\si{\percent}] \\
\midrule
\texttt{muestras26metros} & 20 & 3{,}529 & \textbf{0{,}13} \\
\texttt{muestras34metros} & 21 & 3{,}484 & \textbf{0{,}18} \\
\texttt{muestra18metroslegit} & 8 & 3{,}478 & 0{,}14 \\
\texttt{muestra32} & 21 & 3{,}605 & 0{,}33 \\
\texttt{muestra34} & 33 & 3{,}626 & 1{,}27 \\
\texttt{muestra48} & 24 & 3{,}555 & 2{,}55 \\
\bottomrule
\end{tabular}
\end{center}

Veintiuna sesiones se agolpan entre \SI{3.43}{} y \SI{3.63}{\volt}. No es que el
operador golpeara con precisi\'on de una parte en mil: \textbf{es el techo de la
cadena de adquisici\'on}. La amplitud de esos 297 golpes ---el \SI{46}{\percent}
del conjunto--- \textbf{no es dato}, es un l\'imite del instrumento.}

\seleccionado{\textbf{Hallazgo 3: la repetibilidad real de la fuente, por fin
medida.} La configuraci\'on B (349 golpes, sin saturar, pico mediano
\SI{2.01}{\volt}) da los n\'umeros v\'alidos:

\begin{center}
\begin{tabular}{@{}lrr@{}}
\toprule
M\'etrica & CV mediano & Rango entre posiciones \\
\midrule
Amplitud pico & \SI{8.7}{\percent} & \SIrange{3.8}{14.9}{\percent} \\
Ancho del pulso & \SI{13.7}{\percent} & \SIrange{9.5}{102}{\percent} \\
\textbf{Centroide espectral} & \textbf{\SI{4.2}{\percent}} & \SIrange{3.0}{10.9}{\percent} \\
\bottomrule
\end{tabular}
\end{center}

\textbf{Y el resultado es una buena noticia.} El contenido en frecuencia de la
fuente es \emph{el doble de repetible} que su energ\'ia: \SI{4.2}{\percent}
contra \SI{8.7}{\percent}. Eso importa porque el MASW estima velocidad de fase a
partir de \emph{diferencias de fase entre trazas}, que dependen del contenido
espectral de la fuente, no de su amplitud absoluta ---y la amplitud, adem\'as, se
normaliza al apilar. \textbf{La hip\'otesis que sostiene la apertura sint\'etica
resulta ser m\'as s\'olida en la variable que importa que en la que se ve.} El
centroide se mantiene adem\'as estable en torno a \SIrange{100}{110}{\hertz} en
todo el tendido, sin deriva sistem\'atica con la distancia.}

\pendiente{Tres consecuencias accionables:

\begin{enumerate}[itemsep=0.15em]
    \item \textbf{Rehacer el an\'alisis MASW separando las dos configuraciones},
    o al menos verificar que mezclarlas no sesg\'o la curva de dispersi\'on. Hoy
    el conjunto procesado las trata como homog\'eneas.
    \item \textbf{Bajar la ganancia o cambiar el rango del ADC del nodo martillo}
    ---justamente lo que las cuatro configuraciones de rango existen para
    permitir. El nodo HAMMER no necesita la ganancia agresiva que s\'i necesita el
    ge\'ofono, y aqu\'i se ve el costo de no haberlo ajustado.
    \item El ancho del pulso se dispara a \SI{46}{} y \SI{50}{\meter}
    (CV hasta \SI{313}{\percent}): ah\'i la se\~nal es tan d\'ebil que la medida
    de ancho a media altura pasa a estar dominada por ruido. Es un indicador \'util
    de d\'onde termina el alcance \'util del arreglo, y coincide con la ca\'ida de
    SNR ya reportada m\'as all\'a de los \SI{30}{\meter}.
\end{enumerate}}